\section{Taylor Series}
\label{sec:taylor}
We directly follow up from the previous section by `reversing' the starting point in the following way:
\begin{enumerate}
	\item \Cref{sec:power_series} started with convergent power series $\sum_{n=0}^\infty a_n(x-x_0)^n$. We considered these as functions $f(x) = \sum_{n=0}^\infty a_n(x-x_0)^n$ and then investigated the regularity properties (c.f.~\Cref{thm:cty_power},\Cref{cor:diff_power}, and~\Cref{cor:analytic}). In particular, we deduced that $f(x)$ has infinitely many derivatives.
	\item This section starts with a given function $f(x)$ which we assume has infinitely many derivatives. We then study the formal power series generated by $f(x)$ (c.f. below~\Cref{cor:analytic}) which should be given by
	\begin{equation}
		\label{eq:taylor_series}
		\sum_{n=0}^\infty \frac{f^{(n)}(c)}{n!}(x-c)^n.
	\end{equation}
\end{enumerate}
We already mentioned infinite differentiability, but let us be precise below.
\begin{defn}
	\label{defn:inf_diff_taylor}
	Let $I\subset \R$ and $f:I\to \R$. We say that $f$ is \textit{infinitely differentiable on $I$} provided
	\[
	\forall x\in I,\,\forall k\in \N,\, f^{(k)}(x)\text{ exists.}
	\]
	The set of all functions which are infinitely differentiable on $I$ is $C^\infty(I)$.
	
	Let $c\in\R$ and $\delta>0$ and suppose $f\in C^\infty(c-\delta,c+\delta)$. We define the \textit{Taylor series generated by $f$ around $c$} as the power series in~\eqref{eq:taylor_series}. When $c=0$, the Taylor series~\eqref{eq:taylor_series} is called the Maclaurin series.
\end{defn}
Clearly, the Taylor series~\eqref{eq:taylor_series} converges at $x=c$ and both $f(x)$ and its Taylor series agree at this point. However, a major distinction between the present discussion and previously in~\Cref{sec:power_series} is that~\eqref{eq:taylor_series} might not converge at other points $x\neq c$. Moreover, even if~\eqref{eq:taylor_series} converges at a point $x\neq c$, the Taylor series need not equal $f(x)$.
\begin{eg}
	Let $c=0$ and consider the function
	\[
	f(x) = \left\{
	\begin{array}{cl}
		e^{-\frac{1}{x}}, 	&\text{if }x>0 	\\
		0, 	&\text{if }x\le 0
	\end{array}
	\right..
	\]
	We have the following.
	\begin{enumerate}
		\item (Exercise) $f$ is infinitely differentiable on $\R$. In particular,
		\[
		f^{(n)}(0) = 0,\quad \forall n\in\N.
		\]
		\item However, the Maclaurin series (Taylor series with $c=0$) is identically
		\[
		\sum_{n=0}^\infty \frac{f^{(n)}(0)}{n!}x^n \equiv 0.
		\]
		Since $f(x) \neq 0$ for $x>0$, we see a discrepancy between the function and its Taylor series.
	\end{enumerate}
\end{eg}
In order to study how much a function and its Taylor series differ from each other, let us recall (a version of) Taylor's theorem.
\begin{thm}[Taylor's theorem]
	\label{thm:taylor}
	Let $a<c<b$ be three real numbers and $n\in\N$. Let $f$ be n times differentiable on $(a,b)$. Then, for every $x\in(a,b)\setminus\{c\}$, there exists some $y$ between $c$ and $x$ such that
	\[
	f(x) = \sum_{k=0}^{n-1}\frac{f^{(k)}(c)}{k!}(x-c)^k + \frac{f^{(n)}(y)}{n!}(x-c)^n.
	\]
\end{thm}
\begin{proof}
	\textcolor{blue}{Blank unless necessary.}
\end{proof}
Clearly, the sum on the right-hand side in~\Cref{thm:taylor} is the partial sum of the Taylor series. Therefore, a condition to ensure that a function is equal to its Taylor series is that the so-called \textit{remainder} term
\[
R_n(x) = \frac{f^{(n)}(y)}{n!}(x-c)^n
\]
vanishes as $n\to \infty$.
\begin{thm}
	\label{thm:taylor_eq}
	Let $f\in C^\infty((a,b))$ and fix $c\in (a,b)$. Assume that there exists $\delta>0$ and a constant $M\ge 0$ such that
	\[
	|f^{(n)}(x)|\le M^n,\quad \text{for every }x\in (c-\delta,c+\delta)\cap (a,b)\text{ and every }n\in\N.
	\]
	Then, we have
	\[
	f(x) = \sum_{n=0}^\infty \frac{f^{(n)}(c)}{n!}(x-c)^n,\quad \forall x\in (c-\delta,c+\delta)\cap (a,b).
	\]
\end{thm}
\begin{proof}
	Let $x\in (c-\delta,c+\delta)\cap (a,b)$. According to~\Cref{thm:taylor}, for any $n\in\N$, there exists some $y$ between $x$ and $c$ such that
	\[
	f(x) = \sum_{k=0}^{n-1}\frac{f^{(k)}(c)}{k!}(x-c)^k + R_n(x),\quad R_n(x) = \frac{f^{(n)}(y)}{n!}(x-c)^n.
	\]
	By assumption, we can estimate the remainder by
	\[
	\left|R_n(x)\right|\le \frac{(M\delta)^n}{n!},\quad\forall x\in (c-\delta,c+\delta)\cap (a,b).
	\]
	In particular, the remainder vanishes as $n\to \infty$ because (exercise)
	\[
	\lim_{n\to\infty}\frac{(M\delta)^n}{n!} = 0.
	\]
\end{proof}